\chapter*{Введение}                         % Заголовок
\addcontentsline{toc}{chapter}{Введение}    % Добавляем его в оглавление

\newcommand{\actuality}{}
\newcommand{\progress}{}
\newcommand{\aim}{{\textbf\aimTXT}}
\newcommand{\tasks}{\textbf{\tasksTXT}}
\newcommand{\novelty}{\textbf{\noveltyTXT}}
\newcommand{\influence}{\textbf{\influenceTXT}}
\newcommand{\methods}{\textbf{\methodsTXT}}
\newcommand{\defpositions}{\textbf{\defpositionsTXT}}
\newcommand{\reliability}{\textbf{\reliabilityTXT}}
\newcommand{\probation}{\textbf{\probationTXT}}
\newcommand{\contribution}{\textbf{\contributionTXT}}
\newcommand{\publications}{\textbf{\publicationsTXT}}

{\actuality}\textbf{Актуальность темы исследования.}
Молекулярная спектроскопия высокого разрешения является одним из наиболее точных методов получения
информации о строении молекул, их энергетической структуре, внутримолекулярной динамике и
взаимодействии излучения с веществом. Положения и интенсивности колебательно-вращательных линий
используются при решении задач атмосферной оптики, астрофизики, планетологии, газоанализа, а также при
контроле технологических процессов, в которых необходимо надежно определять состав газовой среды.
Современные спектроскопические базы данных и расчетные списки линий требуют не только большого объема
информации, но и внутренней согласованности спектроскопических параметров, поскольку ошибки в центрах
полос, энергиях уровней или интенсивностях линий непосредственно влияют на результаты моделирования
спектров.

Особое значение имеют изотопические модификации молекул. В рамках приближения Борна--Оппенгеймера
изотопозамещение не меняет электронную потенциальную поверхность как функцию декартовых координат
ядер, однако изменяет массы ядер, моменты инерции, нормальные координаты, частоты колебаний и
параметры эффективного гамильтониана. Поэтому исследование изотопологов позволяет получить
дополнительные ограничения на поверхность потенциальной энергии и поверхность дипольного момента.
Такая информация особенно важна для молекул, у которых экспериментальные данные по отдельным
изотопологам ограничены, а прямое определение всех параметров из независимого спектроскопического
анализа невозможно или недостаточно устойчиво.

В настоящей работе рассматриваются две группы объектов. Первая группа включает тетрагидриды элементов
IV группы GeH$_4$ и SiH$_4$ с тетраэдрической симметрией T$_d$. Эти молекулы относятся к сферическим
волчкам и требуют специального учета симметрии при построении колебательного базиса и матричных
элементов эффективного гамильтониана. Для таких объектов принципиально важны аналитические
представления симметризованных колебательных функций, поскольку именно они позволяют связать
экспериментальные центры полос с гармоническими частотами, ангармоническими параметрами, параметрами
тетраэдрического расщепления и параметрами резонансного взаимодействия. Кроме того, для GeH$_4$ и
SiH$_4$ существенен локально-модовый характер растягивающих колебаний, что необходимо учитывать при
выборе физически обоснованного набора параметров.

Вторая группа объектов включает дейтерированные изотопологи сероводорода D$_2{}^{32}$S, HD$^{32}$S и
HD$^{34}$S в области фундаментальной деформационной полосы $\nu_2$. Сероводород и его изотопологи
представляют интерес для атмосферной оптики, планетологии, астрофизики и задач, связанных с серным
циклом. При этом дейтерированные формы H$_2$S являются удобными объектами для проверки теории
изотопозамещения и методов анализа интенсивностей, поскольку замена H на D существенно меняет
вращательную и колебательную структуру спектра. Дополнительная трудность заключается в
водород-дейтериевом обмене в газовом образце: в ходе длительной регистрации спектров состав смеси
изменяется, а парциальные давления отдельных изотопологов не могут быть определены простым прямым
измерением. Поэтому для корректного определения абсолютных интенсивностей необходима спектроскопическая
процедура оценки парциального давления, основанная на изотопических соотношениях для параметров
эффективного дипольного момента.

Таким образом, актуальность диссертационной работы определяется необходимостью получения согласованных
и достаточно точных спектроскопических параметров для изотопологов молекул различной симметрии,
развитием полуэмпирических методов решения обратной спектроскопической задачи и построением расчетных
наборов данных, пригодных для последующего моделирования спектров высокого разрешения.

\textbf{Степень разработанности темы.}
Классические основы теории колебательно-вращательных спектров, эффективных гамильтонианов,
приближения Борна--Оппенгеймера, теории изотопозамещения и методов симметрии разработаны в трудах
Нильсена, Уотсона, Вильсона, Папоушека, Банкера, Макушкина, Тютерева, Жилинского и других авторов
\cite{nielsen,Watson1968,Watson1977,WilsonDeciusCross1955,Papousek1982,BunkerJensen1998,Makushkin1984,Zhilinskii1984}.
Для молекул типа асимметричного волчка хорошо известны эффективные вращательно-колебательные
гамильтонианы и методы отнесения линий по комбинационным разностям. Для молекул типа сферического
волчка ситуация сложнее: необходимо строить базисные функции, преобразующиеся по неприводимым
представлениям точечной группы, и учитывать тетраэдрическое расщепление вырожденных состояний. В
последние годы были получены аналитические выражения для симметризованных колебательных функций и
матричных элементов гамильтониана молекул типа XY$_4$, что открыло возможность более строгого
полуэмпирического анализа полиад высокой энергии \cite{ulen1,ulen2}.

Несмотря на значительный объем опубликованных данных по основным изотопическим модификациям H$_2$S,
GeH$_4$ и SiH$_4$, информация по менее распространенным или дейтерированным изотопологам остается
неполной. Для GeH$_4$ и SiH$_4$ в литературе имеются экспериментальные центры отдельных полос и
результаты расчетов, однако требуется единое описание нескольких изотопологов в рамках одной
спектроскопической модели. Для D$_2$S и HDS до настоящего исследования существовали данные по части
положений линий, но число отнесенных переходов и диапазон квантовых чисел были существенно меньше, а
сведения об абсолютных интенсивностях полосы $\nu_2$ были ограничены. Это определяет необходимость
дальнейшего анализа спектров высокого разрешения, уточнения параметров энергетических уровней и
получения параметров эффективного дипольного момента.

\textbf{Объектами исследования} являются изотопологи молекул GeH$_4$, SiH$_4$ и H$_2$S, а именно
стабильные изотопологи ${}^{M}$GeH$_4$ ($M=76,74,73,72,70$), ${}^{M}$SiH$_4$ ($M=28,29,30$), а также
D$_2{}^{32}$S, HD$^{32}$S и HD$^{34}$S.

\textbf{Предметом исследования} являются центры колебательных полос, колебательно-вращательные
энергии, параметры эффективного гамильтониана, параметры эффективного оператора дипольного момента,
абсолютные интенсивности спектральных линий и изотопические закономерности, связывающие
спектроскопические параметры различных изотопологов одной молекулы.

{\aim} настоящей работы является получение и уточнение спектроскопических параметров, энергетической
структуры и интенсивностей линий ряда изотопологов молекул различной симметрии на основе методов
молекулярной спектроскопии высокого разрешения, эффективных гамильтонианов и теории
изотопозамещения.

Для~достижения поставленной цели необходимо было решить следующие {\tasks}:
\begin{enumerate}[beginpenalty=10000]
  \item Систематизировать теоретические основы, необходимые для анализа спектров исследуемых молекул:
  гамильтонов подход, приближение Борна--Оппенгеймера, теорию изотопозамещения и элементы аппарата
  неприводимых тензорных операторов.
  \item Применить аналитически симметризованные колебательные функции и матричные элементы
  эффективного гамильтониана для решения обратной спектроскопической задачи молекулы GeH$_4$ и ее
  стабильных изотопологов.
  \item Определить набор колебательных спектроскопических параметров для изотопологов SiH$_4$ с учетом
  локально-модового характера растягивающих колебаний и соотношений теории изотопозамещения.
  \item Выполнить отнесение линий и уточнить энергетическую структуру состояния $(010)$ молекулы
  D$_2{}^{32}$S в области фундаментальной полосы $\nu_2$ на основе высокоразрешенных ИК-спектров.
  \item Разработать и применить процедуру оценки парциального давления D$_2$S в газовой смеси по
  изотопическим соотношениям для параметров эффективного дипольного момента и по измеренным
  интенсивностям пробных линий.
  \item Выполнить анализ положений и интенсивностей линий полосы $\nu_2$ молекул HD$^{32}$S и
  HD$^{34}$S, включая оценку парциального давления HDS в экспериментальном образце.
  \item Сформировать согласованные наборы рассчитанных центров полос и списки переходов с
  индивидуальными интенсивностями, пригодные для последующего моделирования спектров высокого
  разрешения и пополнения спектроскопических баз данных.
\end{enumerate}

{\novelty}
\begin{enumerate}[beginpenalty=10000]
  \item Впервые в рамках единой полуэмпирической схемы, основанной на аналитически симметризованных
  колебательных функциях молекул типа XY$_4$, получены согласованные наборы спектроскопических
  параметров для пяти стабильных изотопологов GeH$_4$ и трех стабильных изотопологов SiH$_4$. Эти наборы
  позволяют рассчитывать центры колебательных полос в области полиад с $N\leq4$ и использовать их как
  основу для дальнейшего моделирования спектров высокого разрешения.
  \item Установлено, что учет локально-модовых соотношений для растягивающих колебаний GeH$_4$ и SiH$_4$
  обеспечивает физически согласованный выбор ангармонических параметров и повышает надежность
  экстраполяции энергетической структуры в области, где прямые экспериментальные данные отсутствуют.
  \item Для полосы $\nu_2$ молекулы D$_2{}^{32}$S существенно расширен набор отнесенных переходов:
  проанализировано 1742 перехода с $J_{\max}=30$ и $K_{a,\max}=21$, получено 535 энергетических значений
  состояния $(010)$ и определен набор из 34 параметров эффективного гамильтониана, воспроизводящий
  исходные энергии со среднеквадратичным отклонением $1.31\times10^{-4}$ см$^{-1}$.
  \item Для D$_2{}^{32}$S получены параметры эффективного оператора дипольного момента полосы $\nu_2$,
  воспроизводящие экспериментальные интенсивности со среднеквадратичным отклонением 2.5~\%; показано,
  что учет изотопических соотношений позволяет определить парциальное давление D$_2$S в условиях
  водород-дейтериевого обмена.
  \item Для полосы $\nu_2$ молекул HD$^{32}$S и HD$^{34}$S выполнено расширенное отнесение линий:
  проанализировано 2684 перехода HD$^{32}$S и 1212 переходов HD$^{34}$S, получены 415 и 263 энергии
  состояния $(010)$ соответственно, а также определены параметры эффективного дипольного момента
  HD$^{32}$S, воспроизводящие интенсивности линий со среднеквадратичным отклонением 4.7~\%.
\end{enumerate}

{\influence} полученных результатов определяется тем, что они дополняют теорию и практику
спектроскопии высокого разрешения для молекул различной симметрии. Теоретическая значимость работы
состоит в развитии и применении полуэмпирического подхода к анализу изотопологов молекул типа XY$_4$ и
молекул типа асимметричного волчка, а также в использовании изотопических соотношений для параметров
эффективного дипольного момента при количественном анализе газовых смесей. Практическая значимость
связана с тем, что полученные центры полос, энергии уровней, параметры эффективного гамильтониана,
параметры эффективного дипольного момента и рассчитанные интенсивности могут быть использованы при
моделировании спектров GeH$_4$, SiH$_4$, D$_2$S и HDS, при подготовке расчетных списков линий, при
сопоставлении с данными спектроскопических баз и при решении прикладных задач атмосферной оптики,
астрофизики, планетологии и газоанализа.

{\methods} исследования включают методы молекулярной спектроскопии высокого разрешения,
Фурье-спектроскопии, теории эффективных гамильтонианов, теории групп и неприводимых тензорных
операторов, теории изотопозамещения, метода комбинационных разностей, метода наименьших квадратов,
полуэмпирического восстановления спектроскопических параметров, анализа профилей спектральных линий
с использованием профиля Хартмана--Трана, а также численного моделирования колебательных и
колебательно-вращательных энергетических структур.

{\defpositions}
\begin{enumerate}[beginpenalty=10000]
  \item Использование аналитически симметризованных колебательных функций и матричных элементов
  эффективного гамильтониана для молекул типа XY$_4$ в сочетании с локально-модовыми и изотопическими
  ограничениями обеспечивает устойчивое определение спектроскопических параметров изотопологов GeH$_4$
  и SiH$_4$ и позволяет предсказывать центры колебательных полос в области полиад $N\leq4$ с точностью,
  согласующейся с имеющимися экспериментальными данными.
  \item Расширенное отнесение линий полосы $\nu_2$ дейтерированных изотопологов сероводорода на основе
  метода комбинационных разностей и эффективного гамильтониана Ватсона обеспечивает описание
  энергетической структуры состояния $(010)$ молекул D$_2{}^{32}$S, HD$^{32}$S и HD$^{34}$S с
  погрешностью, сопоставимой с экспериментальной точностью высокоразрешенных ИК-спектров.
  \item Изотопические соотношения для параметров эффективного дипольного момента позволяют определять
  парциальные давления D$_2$S и HDS в газовой смеси, образующейся при водород-дейтериевом обмене, и тем
  самым обеспечивают корректное восстановление абсолютных интенсивностей линий полосы $\nu_2$ и
  построение согласованных списков переходов для D$_2{}^{32}$S, HD$^{32}$S и HD$^{34}$S.
\end{enumerate}

{\reliability} полученных результатов обеспечивается использованием высокоразрешенных экспериментальных
спектров, большим числом отнесенных переходов, применением независимых комбинационных разностей,
взвешенной подгонкой параметров методом наименьших квадратов, малой величиной остаточных отклонений
<<эксперимент--расчет>>, сопоставлением с ранее опубликованными данными, а также взаимной
согласованностью результатов, полученных для разных изотопологов одной молекулы. Достоверность анализа
интенсивностей дополнительно подтверждается контролем насыщения линий, использованием изолированных
пробных переходов при оценке парциального давления и согласованием главных параметров эффективного
дипольного момента с оценками, следующими из теории изотопозамещения.

{\probation} Основные результаты работы обсуждались на научных семинарах по молекулярной
спектроскопии высокого разрешения и отражены в публикациях, посвященных анализу спектров D$_2$S,
HDS, GeH$_4$ и SiH$_4$. % TODO: при финальной подготовке заменить/дополнить точным перечнем конференций и семинаров.

{\contribution} автора состоит в анализе литературных данных по теме исследования, участии в постановке
задач, подготовке и проверке расчетных схем, обработке экспериментальных спектров, отнесении
спектральных линий, выполнении подгонки параметров эффективного гамильтониана и эффективного
дипольного момента, анализе интенсивностей, подготовке таблиц и графического материала, а также в
обсуждении и интерпретации полученных результатов совместно с научным руководителем и соавторами
публикаций.

{\publications} Основные результаты по теме диссертации отражены в публикациях автора по анализу
высокоразрешенных спектров D$_2$S и HDS, а также по исследованию колебательной структуры молекул
GeH$_4$ и SiH$_4$ \cite{d2s_sydow2019,hds_sydow2019,ulen1,ulen2}. % TODO: при наличии файла author.bib можно заменить этот абзац на автоматический подсчет публикаций.
 % Характеристика работы по структуре во введении и в автореферате не отличается

\textbf{Объем и структура работы.} Диссертация состоит из~введения,
\formbytotal{totalchapter}{глав}{ы}{}{},
заключения и
\formbytotal{totalappendix}{приложен}{ия}{ий}{}.
Полный объём диссертации составляет
\formbytotal{TotPages}{страниц}{у}{ы}{}, включая
\formbytotal{totalcount@figure}{рисун}{ок}{ка}{ков} и
\formbytotal{totalcount@table}{таблиц}{у}{ы}{}.
Список литературы содержит
\formbytotal{citenum}{наименован}{ие}{ия}{ий}.
